\scanpage{402}
\mspage{502}

thanks to the canonical inclusions
\[
\tag{26}
\mu_2\hookrightarrow ST_\rho,\qquad
\mu_2\hookrightarrow ST_\sigma,\qquad
\mu_2\hookrightarrow SG
\]
giving rise to canonical homomorphisms already \unsure{considered} elsewhere
\[
\tag{27}
\Sigma=(\mathbb Z/2\mathbb Z)^3\longrightarrow
\Gamma\mathcal M_{\rho,\sigma}^{\circ+}\deleted{\unsure{+}}
\subset\Gamma\mathcal M_{\rho,\sigma}^{\unsure{+|*}}.
\]

I am going to introduce for the group $\mathcal M_{\rho,\sigma}$ notations which paraphrase the notations $\mathcal M_{0,3}$ etc., in the context of group schemes. The invariant subgroup $SG$ of $\mathcal M_{\rho,\sigma}$ will also be denoted $\widetilde G_{\rho,\sigma}^{+}$ (the analogue of $\widehat{\widetilde G}_{0,3}^{+}$), the quotient
\[
\tag{28}
\deleted{\unsure{\mathcal M}}
\mathcal M_{\rho,\sigma}/\widetilde G_{\rho,\sigma}^{+}
\simeq\tau_{\rho,\sigma}^{!}
\]
\unsure{paraphrases} $\tau_{0,3}^{!}$, we have moreover
\[
\tag{29}
\tau_{\rho,\sigma}^{!}\simeq S\mathcal G_{\rho,\sigma}^{*}(0)
\hookrightarrow T_0\times T_\rho\times N'_\sigma
\]
\margin{\unsure{Seen hanging around a bit} the sign $!$, (\unsure{which} \unsure{has fallen from the sky} --- \unsure{even in the section following the} \uline{\unsure{X}}) it is $\Gamma_{\rho,\sigma}$ (and not $\tau_{\rho,\sigma}$ \unsure{without the exponent} !) that one ought to \unsure{denote} this group by !}

--- in what follows the simpler and more suggestive notation $\tau_{\rho,\sigma}$ will be used in preference to $S\mathcal G_{\rho,\sigma}^{*}(0)$. We thus have
\[
\tag{30}
\diagram{\begin{tikzcd}[ampersand replacement=\&,column sep=small]
1\arrow[r]\&\widetilde G_{\rho,\sigma}^{+}\arrow[r]\arrow[d,phantom,"\rotatebox{90}{$\simeq$}"]\&\mathcal M_{\rho,\sigma}\arrow[r]\&\tau_{\rho,\sigma}^{!}\arrow[r]\&1\\
\&SG\simeq\mathrm{Sl}(2)_S\&\&\&
\end{tikzcd}}
\]
and a canonical splitting as a semidirect product
\[
\tag{31}
\mathcal M_{\rho,\sigma}\simeq\mathcal M_{\rho,\sigma}(0)\mathbin{.}\widetilde G_{\rho,\sigma}^{+}
\]
\margin{401}

\scanpage{403}

where
\[
\tag{32}
\diagram{\begin{tikzcd}[ampersand replacement=\&]
\mathcal M_{\rho,\sigma}(0)\arrow[r,hook]\arrow[dr,"\sim"']\&\mathcal M_{\rho,\sigma}\arrow[d,"\mathrm{can}"]\\
\&\tau_{\rho,\sigma}^{!}
\end{tikzcd}}
\]
is a subgroup isomorphic to $\tau_{\rho,\sigma}^{!}$ --- which corresponds to the subgroup $\mathcal M_{0,3}^{!\sim}(0)$ in $\mathcal M_{0,3}^{\sim}$ ($\S46$, $\uline{\mathrm V}$).

\uline{Remark.} One would like to introduce a $\tau_{\rho,\sigma}$ extension \inserted{(or product)} of $\tau_{\rho,\sigma}^{!}$ by a group such as $\widetilde G_3$ (which would be denoted $\tau_{\rho,\sigma}^{+}$), in relation to a subgroup $\pi_{\rho,\sigma}$ of $\widetilde G_{\rho,\sigma}^{+}$ \deleted{which} such that $\widetilde G_{\rho,\sigma}^{+}/\pi_{\rho,\sigma}\simeq\widetilde G_3$ --- but this is not possible in the context of group \uline{schemes} (even in the modular special case), since $\widetilde G_{\rho,\sigma}^{+}\simeq\mathrm{Sl}(2)_S$ is \deleted{\illegible{1}} \unsure{relatively} connected over $S$. On the other hand one can, in the ``modular case'' $\rho=\begin{pmatrix}0&1\\-1&1\end{pmatrix}$, $\sigma=\begin{pmatrix}0&1\\-1&0\end{pmatrix}$ introduce a canonical section
\[
\tag{33}
\unsure{\dot\tau|\dot t}=\unsure{\dot\tau_\infty|\dot t_\infty}\in\Gamma(\tau_{\rho,\sigma}^{!})
\]
whose representative $\tau_\infty$ in $\Gamma\unsure{\mathcal M}_{\rho,\sigma}^{!}(0)$ is explicitly given as
\[
\tag{34}
\tau_\infty=\begin{pmatrix}-1&0\\0&+1\end{pmatrix}.
\]

(It would be interesting to see whether, outside the
\margin{402}

\scanpage{404}
\mspage{503}

modular case, one can introduce ``anti-involutions'' which generalize the geometric reflections
\[
\tag{35}
\left\{
\begin{aligned}
\deleted{\unsure{\tau}}r_0&=\tau'_\infty=\begin{pmatrix}0&1\\1&\unsure{0}\end{pmatrix}\\
r_1&=-\tau'_0=\begin{pmatrix}-1&1\\0&1\end{pmatrix}\\
r_\infty&=\tau_\infty=\begin{pmatrix}-1&0\\0&+1\end{pmatrix}
\end{aligned}
\right.
\qquad\text{(cf. p. $335'$, $\S42$, $\uline{\mathrm{IV}}$)}
\]
I am convinced that one can. Besides, had we not brought to light in $\mathcal M_{\rho,\sigma}$ a subgroup $\mu_2\times\mu_2\times\mu_2$, \deleted{\unsure{corresponding} \unsure{in}} \inserted{giving} the three remarkable sections (corresponding respective\textsuperscript{ly} to the \deleted{three} sections $-1$ of the three factors $\mu_2$)? But I am getting things mixed up here --- these sections were sections of $\mathcal M_{\rho,\sigma}^{\circ}$ --- whereas the sections \deleted{\unsure{modular}} \inserted{$\tau_\infty,\tau'_\infty,\tau'_0$} \inserted{are of} multiplier
\deleted{in fact the center of $\mathcal M_{\rho,\sigma}$ is \unsure{isomorphic to}}
\inserted{\uline{$Z_{\rho,\sigma}\overset{\mathrm{df}}=(\mathrm{Centr}\,\mathcal M_{\rho,\sigma}^{\circ})$}.}
$-1$, thus correspond to $\varepsilon=\varepsilon'=-1$, thus correspond to the ``least trivial'' section \inserted{$(-1,-1)$} of $\mathcal M_{\rho,\sigma}/\mathcal M_{\rho,\sigma}^{\circ}=\mu_2\times\mu_2$.

\deleted{Let us then return to}
\[
\deleted{(36)}\qquad
\diagram{\begin{tikzcd}[ampersand replacement=\&,column sep=small]
\deleted{\tau_{\rho,\sigma}^{!}\simeq\mathcal G_{\rho,\sigma}^{*}(0)}\arrow[r,hook]\arrow[d,phantom,"\deleted{\rotatebox{90}{$\simeq$}}"]\&\deleted{T_0\times T_\rho\times N'_\sigma}\\
\deleted{\left\{\{u,\underline a,\underline b\}\mid\det\underline u=\det\underline b=\varepsilon'\det\underline a,\ {\varepsilon'}^2=1\right\}}\&
\end{tikzcd}}
\]
\deleted{and let us recall that we have}
\[
\diagram{\begin{tikzcd}[ampersand replacement=\&]
\deleted{T_0}\arrow[r,"\deleted{\sim}"',"\deleted{\delta_0:\ u\mapsto\det u}"]\&\deleted{\mathbb G_m}
\end{tikzcd}}
\]
\deleted{and canonical \unsure{``injective''} homomorphisms}
\[
\diagram{\begin{tikzcd}[ampersand replacement=\&]
\deleted{\mathbb G_m}\arrow[r,"\deleted{j_\rho}"]\&\deleted{T_\rho}\\
\deleted{\mathbb G_m}\arrow[r,"\deleted{j_\sigma}"]\&\deleted{T_\sigma}
\end{tikzcd}}
\]
\deleted{whose \unsure{composite} with $\delta_\rho,\delta_\sigma$ (the determinant homomorphisms) are $\lambda\longrightarrow\lambda^2$. Thus in the second \unsure{side} of $(36)$ there is a}
\margin{403}

\scanpage{405}
\deleted{subgroup of \unsure{trivial centrality}}
\deleted{$\unsure{\mathbb G_m\times}$}
\[
\diagram{\begin{tikzcd}[ampersand replacement=\&]
\deleted{\mathbb G_m\times\mathbb G_m\times\mathbb G_m}\arrow[r,hook]\&\deleted{T_0\times\unsure{T_\rho}\times N'_\sigma}
\end{tikzcd}}
\]

Let us do the calculations to determine the invariants $\alpha\ \beta\ \ldots$ corresponding to $r_0,r_1,r_\infty$.

i) $\displaystyle r_\infty=\tau_\infty=\begin{pmatrix}-1&0\\0&1\end{pmatrix}\in\mathrm{Gl}(2,\mathbb Z)\subset\mathrm{Gl}(2,\widehat{\mathbb Z})^{\pm1}$

We even have: $\tau_\infty\in\unsure{T'_0|B'_0}(\mathbb Z)$, \deleted{so it corresponds to} \inserted{it is the only element} \deleted{$(\alpha,\beta,f)$ with $f=1$,} that is nontrivial in $T'_0(\mathbb Z)\simeq\mathbb G_m(\mathbb Z)\simeq\{\pm1\}$. \deleted{\unsure{The}}

The inner automorphism it defines in $\Pi_{0,3}$ or in $\mathrm{Sl}(2,\mathbb Z)$ --- or their profinite completions --- stabilizes the subgroup $\widetilde L_0$ generated by $\varepsilon_0$ \inserted{(\unsure{i.e.} $L_0$ generated by $\varepsilon_0^2=\ell_0$)}
\[
\tau_\infty(\varepsilon_0)=\varepsilon_0^{-1},\qquad\tau_\infty(\ell_0)=\ell_0^{-1}
\]
$\unsure{(\ast)}$

thus
\[
\tau_\infty\in\mathcal M_{0,3}^{!\sim}[0]
\]
\deleted{\unsure{etc.}} \deleted{\unsure{we find}} it is of the form $u_{\alpha,\beta}$, \deleted{\unsure{we}} the calculations already done give
\[
\alpha=\tau_\infty,\qquad
\unsure{\zeta|\xi}=g=\alpha\rho(\alpha)^{-1}=-\ell_1^{-1}
=-\begin{pmatrix}1&0\\-2&1\end{pmatrix}
\]
(watch the sign $-$ when doing the calculations in $\mathrm{Sl}(2,\mathbb Z)$ and not in $\mathrm{Sl}(2,\mathbb Z)'$ !)
\[
\begin{aligned}
\beta&=\underline1,\qquad\eta=h=1,\\
\mu&=-1,\qquad\varepsilon=\varepsilon_2(\mu)=-1,\qquad
\varepsilon'=\unsure{\varepsilon_3(\mu)|\varepsilon_3^b(\mu)}=-1.
\end{aligned}
\]

we have here $B=0,D=-1$, thus
\[
u=\begin{pmatrix}
\deleted{\unsure{\mu}}\frac1D=\mu&\frac BD=\unsure{\mu D|\mu B}\\
0&1
\end{pmatrix}
=\begin{pmatrix}-1&0\\0&1\end{pmatrix}=\tau_\infty
\]
\[
\underline a=\alpha^{-1}u=\tau_\infty^{-1}\tau_\infty=\underline1,
\qquad\underline b=\beta^{-1}u=\tau_\infty.
\]
thus
\[
\tag{36}
r_\infty=\tau_\infty=
\deleted{\unsure{(\underline u_\infty=-1,\underline a_\infty=1,\underline b_\infty=\tau_\infty)}}
\inserted{(\underline u_\infty=\tau_\infty,\underline a_\infty=1,\underline b_\infty=\tau_\infty)}
\in T_0\times T_\rho\times N'_\sigma.
\]
\inserted{the nontrivial element of ${}_2T_0(\mathbb Z)\simeq{}_2T'_0(\mathbb Z)$}

\uline{N.B.} Let more generally
\[
\underline u=\begin{pmatrix}\mu&\lambda\\0&1\end{pmatrix}
\in\deleted{\unsure{\mathrm{Gl}}}\,B'_0(\widehat{\mathbb Z})\subset\mathrm{Gl}(2,\mathbb Z)
\]
and let us consider the \deleted{inner} automorphism of $\mathrm{Sl}(2,\mathbb Z)'\unsure{\simeq}$ which it induces, hence an automorphism \inserted{($\in\mathcal M_{0,3}^{!\sim}$)} of $\widehat{\widetilde G}_{0,3}^{+}$, what are these invariants $\alpha,\beta$ etc.? We shall be able obvious-
\margin{404}